\ifdefined\MyWebBook\else
\documentclass[../textbook.tex]{subfiles}
\begin{document}
\fi
\bookchapter{条件视觉生成}{ch:conditional-vision}

条件生成要求模型在多样的视觉分布中满足指定信息：文本描述、参考图像、空间布局或时间条件。它既涉及概率路径，也涉及表示压缩、网络结构和求解器。\bookxref{ch:diffusion-probability}推导了加噪与反向过程；本章解释条件如何进入网络、引导如何改变预测，以及潜空间和视频扩展如何改变成本与误差。

离散视觉词元还可以采用自回归因果模型逐项生成，再由相应离散解码器恢复图像，其接口与多模态表征见\bookxref{ch:multimodal-alignment}。

\section{条件生成的形式化}

本章主要符号见表\ref{tab:conditional-vision-symbols}。
\begin{table}[htbp]
\centering
\caption{条件视觉生成主要符号}\label{tab:conditional-vision-symbols}
\begin{tabular}{ll}
\toprule
符号 & 含义\\
\midrule
$\vect{x},\vect{z},c$ & 像素样本、潜变量和条件表示。\\
$H,W,C_x,f,C_z$ & 图像高宽、像素通道数、空间压缩因子与潜通道数。\\
$N,M,d,d_c$ & 视觉位置数、条件位置数、视觉与条件特征宽度。\\
$a_t,b_t$ & 扰动路径中的信号系数与噪声系数。\\
$\vect{s}_t,\vect{\epsilon}_{\vect{\theta}},w$ & 得分函数、噪声预测与本章 CFG 引导系数。\\
$\tau,\vect{u}_{\vect{\theta}}$ & 从噪声走向数据的连续时间与速度场。\\
$F,F',p_t,p_h,p_w$ & 原始及潜空间帧数与时空补丁尺寸。\\
\bottomrule
\end{tabular}
\end{table}

文本生成图像关注条件遵循与分布质量，图像编辑还要求保留未编辑区域，视频生成进一步要求跨帧身份、运动与时间一致性。条件可以作为序列、空间张量或全局向量进入网络，转换成一段文本只是其中一条路径。训练时必须明确每种条件的缺失表示、对齐关系和可见范围。

\section{去噪网络}

\subsection{U-Net 的多尺度路径}
\term{U形网络}{U-Net}{}由下采样路径、瓶颈与上采样路径构成。下采样减少空间位置数，扩大后续局部算子的有效感受野；跳跃连接将较高分辨率特征传给对应上采样层。全局建模能力取决于层数、卷积感受野与注意力配置。

时间编码可采用
\begin{equation}
 e_{2j}(t)=\sin(\omega_j t),\qquad e_{2j+1}(t)=\cos(\omega_j t),
\end{equation}
再经多层感知机形成调制向量。对特征 $\vect{h}$，常见条件调制为
\begin{equation}
 \vect{h}'=\vect{\gamma}(\vect{e}_t,c)\odot\operatorname{Norm}(\vect{h})+\vect{\beta}(\vect{e}_t,c).
\end{equation}
$\gamma,\beta$ 沿空间维广播；其形状必须与通道一致。离散步编号、归一化时间与对数信噪比是三种不同的输入，替换时间编码时必须保持训练契约。

\begin{figure}[htbp]
\centering
\includegraphics[width=.92\linewidth]{\subfix{../../figures/theory/conditional-unet-flow.pdf}}
\caption[条件 U-Net 结构与跨层跳跃连接]{条件 U-Net 结构与跨层跳跃连接。时间与语义条件嵌入各尺度残差块，跳跃连接传递多尺度空间特征且仍参与上采样融合计算。}\label{fig:vision-unet}
\end{figure}

图\ref{fig:vision-unet}中的时间与语义条件在不同尺度影响预测，跳跃特征仍要经过上采样分支的可训练计算。

\subsection{交叉注意力和联合注意力}
视觉特征 $\mat{X}\in\mathbb R^{N\times d}$ 和条件特征 $\mat{C}\in\mathbb R^{M\times d_c}$ 通过
\begin{equation}
 \begin{aligned}
  \mat{Q} &= \mat{X}\mat{W}_Q,\quad \mat{K}=\mat{C}\mat{W}_K,\quad \mat{V}=\mat{C}\mat{W}_V, \\
  \mat{A} &= \operatorname{softmax}\left(\frac{\mat{Q}\mat{K}^\top}{\sqrt{d_h}}+\mat{B}\right)\mat{V}.
 \end{aligned}
\end{equation}
建立交叉注意力。这里 $\mat{W}_Q\in\mathbb R^{d\times d_h}$，$\mat{W}_K\in\mathbb R^{d_c\times d_h}$，$\mat{W}_V\in\mathbb R^{d_c\times d_v}$，输出 $\mat{A}\in\mathbb R^{N\times d_v}$；$\mat{B}$ 排除条件填充等不可见位置。单头打分成本随 $NMd_h$ 增长，视觉自注意力则是 $N^2d_h$。

把视觉与文本映射到兼容宽度后联合处理，可以让两种表示相互更新，但会产生约 $(N+M)^2$ 的注意力配对数。是否双向更新、是否使用独立投影和调制，是具体架构选择。文本编码器、分词器、截断规则和空条件表示都属于制品，替换编码器名称远不足以对齐两者。

\subsection{Transformer 主干与生成目标}
\term{扩散 Transformer}{Diffusion Transformer}{DiT}把潜变量切分为补丁并使用 Transformer 预测去噪目标\cite{peebles2023dit}。若潜图为 $H_z\times W_z\times C_z$、补丁边长为 $p$，则位置数为 $N=\frac{H_zW_z}{p^2}$，每补丁输入宽度为 $p^2C_z$。输出投影恢复每补丁所需通道，再重排为潜图。

时间或类别可调制归一化与残差门控。补丁变小会增加位置数；边长减半使 $N$ 增至4倍，完整自注意力配对数增至16倍。主干结构与训练目标可分别选择：噪声预测、速度参数化和路径速度回归都可以由不同主干实现。

\section{无分类器引导}

\subsection{噪声外推}
\term{得分函数}{Score Function}{}在本节专指连续概率密度对空间输入的对数梯度 $\vect{s}_t(\vect{x})=\nabla_{\vect{x}}\log p_t(\vect{x})$，区别于分类任务中的对数几率输出。给定条件 $c$，贝叶斯关系给出
\begin{equation}
 \nabla_{\vect{x}}\log p_t(c\mid \vect{x})=\vect{s}_t(\vect{x}\mid c)-\vect{s}_t(\vect{x}).
\end{equation}
将条件信息梯度乘以 $w$ 加到无条件得分上，得到
\begin{equation}
 \widehat s_t=(1-w)\vect{s}_t(\vect{x})+w \vect{s}_t(\vect{x}\mid c).
 \label{eq:vision-cfg-score}
\end{equation}
对同一高斯扰动 $\vect{x}_t=a_t\vect{x}_0+b_t\vect{\epsilon}$，密度可微且交换积分与微分合法时，
\begin{align}
 \nabla_{\vect{x}_t}\log p_t(\vect{x}_t\mid c)
 &=\mathbb E\left[-\frac{\vect{x}_t-a_t\vect{x}_0}{b_t^2}\,\middle|\,\vect{x}_t,c\right]\\
 &=-\frac{\mathbb E[\vect{\epsilon}\mid \vect{x}_t,c]}{b_t}.
\end{align}
因此噪声回归的最优条件均值与得分成比例。在相同 $b_t>0$ 下，式\eqref{eq:vision-cfg-score}对应
\begin{equation}
 \widehat{\vect{\epsilon}}=\vect{\epsilon}_{\vect{\theta}}(\vect{x}_t,t,\varnothing)
 +w\left[\vect{\epsilon}_{\vect{\theta}}(\vect{x}_t,t,c)-\vect{\epsilon}_{\vect{\theta}}(\vect{x}_t,t,\varnothing)\right].
 \label{eq:vision-cfg-noise}
\end{equation}
\term{无分类器引导}{Classifier-Free Guidance}{CFG}通过条件丢弃训练，让生成模型提供两种估计，无须另训分类器\cite{ho2022cfg}。丢弃的对象是语义条件；空条件的编码方式也必须在训练中确定。

\subsection{外推精度边界}
本章约定 $w=0$ 对应无条件预测，$w=1$ 对应无引导的条件预测，$w>1$ 实施外推引导。部分文献记为 $(1+s)\vect{\epsilon}_c-s\vect{\epsilon}_u$，其标量满足 $s=w-1$；跨库比对引导强度时须按公式换算。

在固定时间且归一化积分存在时，式\eqref{eq:vision-cfg-score}是形式密度 $p_t(\vect{x}\mid c)^w p_t(\vect{x})^{1-w}$ 的对数梯度。但不同时间的这种密度族未必满足原前向扩散关系，据此断言最终采样分布恰为 $p_0(\vect{x}\mid c)^w p_0(\vect{x})^{1-w}$ 缺少依据。模型近似与求解误差还会进一步改变结果。

设无条件与条件预测分别为 $(0.2,-0.1)$ 与 $(0.5,0.3)$，则 $w=2$ 得 $(0.8,0.7)$。引导放大了条件与无条件预测之间的差向量。若两支预测误差为 $\vect{e}_u,\vect{e}_c$，引导误差满足
\begin{equation}
 \|\vect{e}_{\rm cfg}\|\le |1-w|\|\vect{e}_u\|+|w|\|\vect{e}_c\|.
\end{equation}
较大的引导强度可能放大误差、改变多样性并产生伪影。以负面提示替代空条件时，基线对应负面提示下的条件预测，引导方向因而取决于正、负两个条件分支。\footnote{将两支合并为一个批次可以减少某些调度开销，但仍需处理两组条件前向，成本仍高于只计算一支。}

\section{潜空间扩散}

\subsection{变分自编码器}
\term{变分自编码器}{Variational Autoencoder}{VAE}以 $q_{\vect{\phi}}(\vect{z}\mid \vect{x})$ 近似后验，解码器为 $p_{\vect{\psi}}(\vect{x}\mid \vect{z})$。由 Jensen 不等式，
\begin{equation}
 \log p_{\vect{\psi}}(\vect{x})\ge
 \mathbb E_{q_{\vect{\phi}}}\log p_{\vect{\psi}}(\vect{x}\mid \vect{z})
 -D_{\rm KL}(q_{\vect{\phi}}(\vect{z}\mid \vect{x})\|p(\vect{z})).
\end{equation}
对对角高斯 $q_{\vect{\phi}}=\mathcal N(\mu,\operatorname{diag}\sigma^2)$ 与标准正态先验，KL 为
\begin{equation}
 \frac12\sum_j(\mu_j^2+\sigma_j^2-\log\sigma_j^2-1).
\end{equation}
重参数化 $\vect{z}=\mu_{\vect{\phi}}(\vect{x})+\sigma_{\vect{\phi}}(\vect{x})\odot\eta$、$\eta\sim\mathcal N(0,I)$ 将随机输入与参数分离，使梯度沿均值与尺度传播。实际视觉自编码器可能加入感知或对抗重建目标，其总目标由各重建项与正则项共同定义。

\term{潜空间扩散}{Latent Diffusion}{}先学习压缩表示，再在该表示上训练生成过程\cite{rombach2022ldm}。图像为 $H\times W\times C_x$、压缩率为 $f$ 时，潜元素数为 $\frac{HWC_z}{f^2}$，原图与潜空间元素比为 $\frac{f^2C_x}{C_z}$。例如 $512\times512\times3$ 到 $64\times64\times4$，比值为48，空间位置数减少64倍。端到端加速比还取决于编码、解码以及生成主干的通道宽度和计算量。

\subsection{缩放协议与不可恢复信息}
若生成器训练使用 $\vect{z}'=s\vect{z}$，采样结束后必须先得到 $\vect{z}=\frac{\vect{z}'}{s}$ 再解码。遗漏缩放改变解码器输入分布；缩放值使用对应制品的配置。像素归一化、色域、裁剪、潜通道、后验采样或均值选择也必须一致。

编码器丢失的信息留在表示瓶颈处，后续生成器再大也补不回来。细小文字、线条与局部身份细节尤其需要独立验证。自编码器重建误差和生成误差属于不同来源：先对真实输入做编码解码对照，才能判断瓶颈是否已出现在表示层。

\section{流匹配}
\optional
\subsection{条件速度回归}
本节以 $\tau\in[0,1]$ 从噪声走向数据，端点记为 $\vect{Z}_0$ 与 $\vect{Z}_1$，区别于\bookxref{ch:diffusion-probability}从干净样本出发的离散加噪记号。取一个端点联合分布，定义线性路径
\begin{equation}
 \begin{aligned}
  \vect{Z}_\tau &= (1-\tau)\vect{Z}_0+\tau \vect{Z}_1, \\
  \vect{U} &= \vect{Z}_1-\vect{Z}_0.
 \end{aligned}
\end{equation}
\term{流匹配}{Flow Matching}{}可以通过条件路径监督速度场\cite{lipman2023flowmatching}。平方损失
\begin{equation}
 \mathcal L=\mathbb E\|\vect{u}_{\vect{\theta}}(\vect{Z}_\tau,\tau,c)-\vect{U}\|^2
\end{equation}
在固定 $(\vect{z},\tau,c)$ 下的最优解为 $\vect{u}^*(\vect{z},\tau,c)=\mathbb E[\vect{U}\mid \vect{Z}_\tau=\vect{z},c]$。证明来自平方展开：任意预测与条件均值之差的平方，加上与预测无关的条件方差，构成全部条件风险。

对光滑紧支撑测试函数 $\varphi$，若微分与期望可以交换，
\begin{equation}
 \frac{d}{d\tau}\mathbb E[\varphi(\vect{Z}_\tau)]
 =\mathbb E[\nabla\varphi(\vect{Z}_\tau)^\top \vect{U}]
 =\mathbb E[\nabla\varphi(\vect{Z}_\tau)^\top \vect{u}^*(\vect{Z}_\tau,\tau,c)].
\end{equation}
因此该边缘速度满足对应密度路径的弱形式连续性方程。若速度场具备适当正则性与可积性，求解 $\frac{d\vect{Z}_\tau}{d\tau}=\vect{u}^*(\vect{Z}_\tau,\tau,c)$ 可以运输这些边缘分布。端点配对可以独立，也可来自其他耦合；它会改变条件速度与学习难度。

\subsection{直线路径的生成边界}
\optional

\begin{remark}[样本路径与生成轨迹]
训练样本路径是直线，但同一中间位置可能对应多对端点，其条件平均速度会随位置与时间变化。边缘速度场的积分轨迹通常偏离训练样本的直线，生成仍需按相应求解器积分。
\end{remark}

Euler 更新为 $z_{k+1}=z_k+h_k u_\theta(z_k,\tau_k,c)$。若场足够光滑，单步局部截断误差为二阶、固定区间的全局误差通常为一阶；网络误差另计。
\begin{example}[Euler 离散化的误差]
以一维 $u(z)=z,z(0)=1$ 为例，精确终点为 $\conste$，单步 Euler 得2，两步得2.25，都不精确。求解步数决定该离散化的精度。
\end{example}

扩散中的 $v$ 参数化常表示\emph{信号与噪声的特定线性组合}，流匹配的 $u$ 则是\emph{路径导数}。接口中的 velocity 应对应各自定义的预测量。预测对象、时间方向和求解器必须共同确定。现代条件视觉生成的主流技术路线与架构对比见表\ref{tab:vision-generation-taxonomy}。

\begin{table}[htbp]
\centering
\small
\setlength{\tabcolsep}{3pt}
\caption{现代条件视觉生成主流范式对比}\label{tab:vision-generation-taxonomy}
\begin{tabularx}{\linewidth}{@{}p{1.8cm}p{2.0cm}p{2.6cm}p{1.8cm}X@{}}
\toprule
生成范式 & 空间与主干 & 概率目标与过程 & 代表系统 & 核心优势与工程瓶颈 \\
\midrule
像素扩散 & 级联 U-Net & 离散扰动链噪声预测 & DDPM, Imagen & 细节保真无自编码损耗；显存计算开销巨大 \\
潜空间扩散 & 交叉注意 U-Net & 连续潜空间噪声/$v$ 预测 & SD 1.5/XL & 空间压缩 $f=8$ 降低负载；上限受 VAE 解码约束 \\
扩散 DiT & 纯 Transformer & 连续时间 adaLN 调制 & PixArt, Sora & 扩展规律显著；密集潜补丁下注意力计算量大 \\
流匹配 (FM) & MMDiT / DiT & 最优传输条件速度场回归 & SD3, FLUX.1 & 采样路径平直仅需 10--20 步；需大容量主干拟合 \\
离散自回归 & 因果解码器 & 离散码本自回归最大似然 & DALL-E, Parti & 与 LLM 范式统一；逐词元采样串行延迟偏高 \\
\bottomrule
\end{tabularx}
\end{table}

\section{条件生成的应用}
\optional
图像编辑将参考图、遮罩或结构控制作为额外条件。若要求未编辑区域严格保持，应区分像素级约束与模型倾向；潜空间遮罩经过编码解码可能影响边缘及邻近区域。多次去噪中的已知区域替换还需匹配当前噪声尺度；把干净像素与高噪状态混合，得到的已不是原模型所定义的状态。

视频潜变量可写为 $[B,C_z,F',H',W']$，时空补丁数为
\begin{equation}
 N=\frac{F'}{p_t}\frac{H'}{p_h}\frac{W'}{p_w},
\end{equation}
这里假设尺寸可整除，真实实现的补齐与边界另计。设每帧有 $S$ 个空间位置，共 $F'$ 帧，全时空注意力配对数为 $F'^2S^2$；空间与时间分解注意力约为 $F'S^2+SF'^2$。分解注意力通过时空正交分解降低了计算复杂度，但牺牲了任意跨帧补丁间的直接单层交互通路。

\optional
\begin{example}[视频潜序列的注意力规模]
例如16个潜帧，每帧 $32\times32$ 位置，补丁为 $1\times2\times2$，则 $S=256,N=4096$，完整打分矩阵每头有约1678万项。帧数翻倍使全注意力项增至4倍，而空间分支仅翻倍。帧率、时长与潜时间压缩共同决定 $F'$；某些因果时间编码器对首帧采用特殊处理，$F'=\frac{F}{f_t}$ 只适用于不含这类处理的实现。
\end{example}

音视频联合生成需要共同时间基准、采样率与对齐条件。声画同步通过时间对应关系检验，物理预测通过状态变化检验。世界模型还需覆盖动作条件与状态空间，并在闭环任务中评价控制表现。

\section{生成模型的执行一致性}

\begin{algorithm}[条件潜空间生成的数据流]

\begin{enumerate}
\item 固定像素预处理、编码器与潜缩放，取得训练潜变量及条件表示；记录条件缺失规则。
\item 按训练目标采样时间和噪声，构造扰动或路径状态，计算噪声、样本或路径速度目标。
\item 训练主干，确保输出形状、监督对象与时间权重一致；条件丢弃只按已定义策略执行。
\item 推理时从匹配先验初始化，用匹配的预测参数化和时间方向逐步求解；需要引导时组合两支预测。
\item 恢复潜缩放并解码，保留精确组件版本、生成配置与输出来源记录。
\item 分别评价重建限制、条件遵循、质量、多样性、编辑保持与视频时间一致性；不以单张图代替分布评估。
\end{enumerate}
\end{algorithm}

\begin{figure}[htbp]
\centering
\includegraphics[width=.92\linewidth]{\subfix{../../figures/theory/conditional-generation-loop.pdf}}
\caption[条件视觉迭代生成与解码流程]{条件视觉迭代生成与解码流程。潜在空间去噪网络与数值求解器迭代协作更新状态轨迹，像素解码仅在去噪终点单次执行。}\label{fig:vision-generation}
\end{figure}

图\ref{fig:vision-generation}显示预测与状态更新的分工；只有组件协议相容时，循环才具有所声明的概率含义。质量评估需同时检查条件遵循、分布多样性及任务细节：图文相似度评价表示空间中的匹配，分布距离评价总体样本特征；计数、否定、局部文字与条件遵循通过专门任务评价。

服务预算应同时限制分辨率、帧数、采样步数、候选数和并发；这些量共同决定请求的工作量。组件清单包含主干、文本编码器、分词器、自编码器、缩放、求解器与控制模块。部署前应逐项核对所用版本的代码、权重获取方式和许可用途，并保留对应记录。


% BEGIN GENERATED INTERVIEWS

% END GENERATED INTERVIEWS
\chapterreferences
\myendchapterdocument
